\Needspace{18\baselineskip}
\fancyhead[L]{\small KOMET VON 1807.}
\begin{center}
KOMET VON 1807.
\end{center}

\begin{center}
\fontsize{10}{12}\selectfont\begin{tabular}{c}
\hline
Göttingische gelehrte Anzeigen. Stück 32. Seite 313..315. 25. Februar 1808. \\
\hline
\end{tabular}
\end{center}

Das ungünstige Wetter, das den ganzen Winter über anhielt, erlaubte an der hiesigen Sternwarte nur wenige Beobachtungen des letzten Kometen, und diese meist unter ungünstigen Bedingungen. Da die Beobachtungen inzwischen überall zu Ende sein dürften, stellen wir hier alle zusammen, die seit der Ankunft von Prof. \textsc{Gauss} hier möglich waren.


\begin{center}\fontsize{9}{11}\selectfont\setlength{\tabcolsep}{2pt}\renewcommand{\arraystretch}{1.15}\begin{tabular}{lrrrrr}
& & \multicolumn{2}{c}{Scheinbare} & \multicolumn{2}{c}{Unterschied}\\
& Mittlere Zeit & \shortstack{Rektaszension} & \shortstack{nordliche\\Deklination} & \shortstack{Rektasz.} & Deklination\\\hline
1807 Dez. 17. & $11^{\text{u}}\,24^{\text{m}}\,44^{\text{s}}$ & $305^\circ\;54'\;0''$ & $45^\circ\;4'\;30''$ & $+49$ & $-28$\\
29 & $7\;59\;21$ & $318\;57\;41$ & $46\;41\;48$ & $+6$ & $-25$\\
29 & $11\;36\;33$ & $319\;6\;35$ & $46\;43\;11$ & $+49$ & $-55$\\
1808 Jan. 4 & $7\;22\;12$ & $325\;11\;19$ & $47\;12\;46$ & $-18$ & $-21$\\
4 & $9\;13\;57$ & $325\;15\;53$ & & $-6$ & \\
4 & $17\;30\;26$ & $325\;36\;38$ & & $+18$ & \\
31 & $9\;31\;34$ & $349\;46\;43$ & $48\;9\;26$ & $-84$ & $+24$\\
\end{tabular}\end{center}

\par\indent Bis auf die letzte wurden alle diese Beobachtungen mit dem Kreismikrometer durch Vergleich mit Sternen aus der \emph{histoire céleste} angestellt. Bei der letzten diente ein Stern siebter Größe zum Vergleich, dessen Position bis dahin nicht bestimmt war. Erst am 14. Februar ließ sich seine Position beobachten. Die Beobachtung ergab


\begin{center}\fontsize{10}{12}\selectfont\setlength{\tabcolsep}{2pt}\renewcommand{\arraystretch}{1.15}\begin{tabular}{lr}
scheinbare Rektaszension des Sterns .\;.\;.\;.\; & $350^\circ\;14'\;38''$\\
scheinbare nördliche Deklination .\;.\;.\;.\;.\;.\;.\; & $48\;4\;33$\\
\end{tabular}\end{center}

\noindent Die obige Beobachtung wurde anhand dieser Position reduziert, die noch zu berichtigen ist.

\par\indent Die aus der Monatlichen Correspondenz bereits bekannten parabolischen Elemente des Kometen, die Prof. \textsc{Gauss} schon Anfang November berechnet hatte, stimmten mit diesen späteren Beobachtungen folgendermaßen überein [die Unterschiede sind in diesem Abdruck oben neben den Beobachtungen angegeben.]

\Needspace{12\baselineskip}
\par\indent Es wäre kaum der Mühe wert, diese Elemente anhand der Beobachtungen nochmals zu verbessern, zumal schon wegen des schwachen Lichts des Kometen seine Ein- und Austritte am 31. Januar nur mit Mühe beobachtet werden konnten und die letzten Beobachtungen daher keine besonders hohe Genauigkeit erreichen konnten. Wenn man an anderen Sternwarten nicht viel mehr Glück hatte, wird sich trotz der langen Sichtbar\SourceLeaf{304}keit dieses Kometen schwerlich eine entscheidende Aussage über die Abweichung seiner Bahn von einer Parabel treffen lassen. Gewiss kann die Abweichung aber in jedem Fall nur sehr klein sein.

\par\indent Am 14. Februar war der Komet im \textsc{Herschel}-Teleskop gerade noch zu sehen, doch an eine eigentliche Beobachtung war nicht mehr zu denken.
